% from ega3-6-fr.tex
\oldpage[III]{141}

de $Y$, $f^{-1}(Y)$ est un schéma, et pour appliquer ce qui précède, il
suffit de vérifier que pour une suite exacte
\[
  0\longrightarrow\mathscr K'^\bullet
  \longrightarrow\mathscr K^\bullet
  \longrightarrow\mathscr K''^\bullet
  \longrightarrow0
\]
de complexes de $\mathscr O_X$-Modules quasi-cohérents, l'homomorphisme
\[
  \partial:
  \mathscr H^n(f,\mathscr K''^\bullet)|V
  \longrightarrow
  \mathscr H^{n+1}(f,\mathscr K'^\bullet)|V
\]
ne dépend pas du recouvrement ouvert affine $\mathfrak U$ de $f^{-1}(V)$
utilisé pour le définir. Mais cela résulte de ce que, si $\mathfrak U'$ est
un recouvrement ouvert affine plus fin que $\mathfrak U$, le diagramme
\[
  \xymatrix@C=34pt@R=11pt{
    \mathbf H^\bullet(\mathfrak U,\mathscr K^\bullet)
      \ar[rrd]^-{\sim}\ar[dd]^{\iota}
    && \\
    &&
    \mathbf H^\bullet(f^{-1}(V),\mathscr K^\bullet)
    \\
    \mathbf H^\bullet(\mathfrak U',\mathscr K^\bullet)
      \ar[rru]_-{\sim}
    &&
  }
\]
d'isomorphismes canoniques est commutatif, ainsi que le diagramme
\[
  \xymatrix@C=46pt@R=25pt{
    \mathbf H^n(\mathfrak U,\mathscr K''^\bullet)
      \ar[r]^-{\partial}\ar[d]^{\iota}
    &
    \mathbf H^{n+1}(\mathfrak U,\mathscr K'^\bullet)
      \ar[d]^{\iota}
    \\
    \mathbf H^n(\mathfrak U',\mathscr K''^\bullet)
      \ar[r]_-{\partial}
    &
    \mathbf H^{n+1}(\mathfrak U',\mathscr K'^\bullet).
  }
\]

Lorsque l'une des conditions précédentes est remplie et que
$\mathscr K^\bullet\to\mathscr K'^\bullet$ est un homotopisme (6.1.4),
l'isomorphisme correspondant
$\mathscr H^\bullet(f,\mathscr K^\bullet)\xrightarrow{\sim}
\mathscr H^\bullet(f,\mathscr K'^\bullet)$ est alors un isomorphisme de
$\partial$-foncteurs (0, 11.4.4).
\end{env}

\begin{env}[6.2.7]
Tout ce qui précède s'applique naturellement sans changement (sinon de
notations) à un complexe $\mathscr K_\bullet$ de $\mathscr O_X$-Modules
quasi-cohérents dont l'opérateur de dérivation est de degré $-1$; il suffit
de considérer le complexe $\mathscr K^\bullet=(\mathscr K^i)$ où
$\mathscr K^i=\mathscr K_{-i}$ pour tout $i\in\mathbf Z$.
\end{env}

\subsection{Hypertor de deux complexes de modules}
\label{subsection:III.6.3.fr}

\begin{env}[6.3.1]
Soient $A$ un anneau commutatif, $P_\bullet$, $Q_\bullet$ deux complexes de
$A$-modules dont les opérateurs de dérivation sont de degré $-1$; soit
$L_{\bullet\bullet}$ (resp. $M_{\bullet\bullet}$) une résolution projective
de Cartan-Eilenberg de $P_\bullet$ (resp. $Q_\bullet$) (0, 11.6.1);
$L_{\bullet\bullet}\otimes_A M_{\bullet\bullet}$ est alors (pour la somme
des premiers degrés et la somme des seconds degrés) un bicomplexe (à
opérateurs de dérivation de degré $-1$), dont l'homologie
$H_\bullet(L_{\bullet\bullet}\otimes_A M_{\bullet\bullet})$ ne dépend pas des
résolutions de Cartan-Eilenberg $L_{\bullet\bullet}$,
$M_{\bullet\bullet}$ choisies, et est par définition l'hyperhomologie du
bifoncteur $P_\bullet\otimes_A Q_\bullet$ en $P_\bullet$ et $Q_\bullet$
(0, 11.6.5). Nous poserons par définition
\[
  \mathbf{Tor}_n^A(P_\bullet,Q_\bullet)
  =H_n(L_{\bullet\bullet}\otimes_A M_{\bullet\bullet})
  \tag{6.3.1.1}\label{III.6.3.1.1.fr}
\]

